\section{Why the top cup power does not vanish}\label{sec:nonvanishing}

The cup-product calculation has reduced nonvanishing to a question about
functions.  We must show that the product in \eqref{eq:top-cochain} cannot be
the restriction of a coboundary.  Two features of the construction work
together: approaching infinity forces eventual constancy in a coordinate,
and the strictly increasing coordinate cardinalities let us use this
constancy simultaneously for every choice of the other coordinates.

Recall that \(D_i=A_i\times\F\), with \(|A_i|=\aleph_i\).  Write

\[
C_r=\prod_{i=0}^{r}D_i,
\qquad
\sigma_r((a_i,s_i)_{i=0}^{r})=(a_i,s_i+1)_{i=0}^{r}.
\]

Thus \(C_r\) consists of the all-finite configurations, before passing to
the quotient.  The pullback of the top product cochain is

\[
F_r((a_i,s_i)_{i=0}^{r})
=\prod_{j=0}^{r-1}(s_j+s_{j+1}).
\]

We use the empty-product convention \(F_0=1\).  Throughout this section,
all function values and sums are in \(\F\).

The function \(F_r\) depends only on the bits, not on the base points
\(a_i\).  For instance, it takes the value one whenever consecutive bits
are different.  This observation alone says nothing about cohomological
nonvanishing: a nonzero function can still be a coboundary.  The role of the
large base sets is to restrict which invariant functions can occur in a
boundary, through the continuity requirements at the points at infinity.
We exploit those requirements without trying to describe all possible
continuous functions on the intersections.

\subsection{What a boundary would imply}

Call a function \(f:C_r\to\F\) \emph{almost constant in coordinate \(i\)}
if, after fixing every other coordinate, its restriction to \(D_i\) is
constant outside a finite set.  Neither the exceptional set nor the
eventual constant is required to be independent of the fixed coordinates.

\begin{lemma}\label{lem:boundary-necessary}
Suppose \(r\geq1\).  If the top product cochain is a coboundary, then

\[
F_r=\sum_{i=0}^{r}f_i,
\]

where every \(f_i:C_r\to\F\) is invariant under \(\sigma_r\) and almost
constant in coordinate \(i\).
\end{lemma}

\begin{proof}
An ordered cochain one degree below the top has a component on each
intersection

\[
W_{\widehat i}=\bigcap_{j\neq i}W_j.
\]

These components are locally constant \(\F\)-valued functions.
The top differential is their sum after restriction to the full
intersection: its alternating signs disappear in characteristic two.
Pulling the components back to \(C_r\) gives the functions \(f_i\).
Because they are pulled back from the orbit space, they are invariant
under the simultaneous flip.

Fix every coordinate except \(i\).  Allowing coordinate \(i\) to vary
through its entire one-point compactification defines a continuous map

\[
D_i^+\longrightarrow W_{\widehat i}.
\]

This map is defined even at infinity.  Indeed, \(r\geq1\), so at least one
other coordinate remains finite; the removed all-infinite orbit is never
encountered.  Compose with the component function on \(W_{\widehat i}\).
Since \(\F\) is discrete, continuity at infinity says that this composite
equals its value at infinity on a cofinite subset of \(D_i\).  This is
exactly the required almost constancy.  Notice that this concerns the
whole doubled set \(D_i=A_i\times\F\).  Outside a finite set, both bit
values therefore have the same eventual value.  Two unrelated eventual
values on the separate bit-zero and bit-one subsets would not suffice
for the argument below.
\end{proof}

Only this necessary condition is used.  In particular, we do not claim
that an arbitrary family of coordinatewise almost-constant functions
extends to jointly continuous functions on the deleted intersections.
\leanmark{boundary}

\subsection{One choice avoiding all exceptional sets}

The exceptional set just obtained can depend on many other coordinates.
The following elementary counting argument supplies a single choice
that avoids all of them.

\begin{lemma}[Uniform thinning]\label{lem:uniform-thinning}
Let \(P\) be infinite and \(|P|<|A|\).  For each \(y\in P\), let
\(N_y\subseteq A\times\F\) be finite.  There is an \(a\in A\) such that
both \((a,0)\) and \((a,1)\) lie outside every \(N_y\).
\end{lemma}

\begin{proof}
Let \(\pi:A\times\F\to A\) be the first-coordinate projection.  Since
each \(\pi(N_y)\) is finite,

\[
\left|\bigcup_{y\in P}\pi(N_y)\right|
\leq |P|\cdot\aleph_0
=|P|
<|A|.
\]

Choose \(a\) outside this union.  No point of either form \((a,0)\) or
\((a,1)\) can belong to any of the exceptional sets.
\end{proof}

For the application, the index set is \(C_{r-1}\), whose cardinality is
\(\aleph_{r-1}\), while the last base set has cardinality \(\aleph_r\).
This strict inequality is the reason for increasing the coordinate sizes.
The argument uses no uniform bound on the sizes of the finite exceptional
sets, and no common eventual value.  It is simply a bound on a union of
finite sets, so no regularity assumption on the larger cardinal is being
used.

Choosing a point separately for each configuration would be much weaker.
It would give a base point depending on the remaining coordinates, and
then taking the associated slices could destroy almost constancy in
those coordinates.  The uniform choice instead gives two fixed slices
of every function.  Their algebraic and continuity properties can then
be compared directly.  This is the only place where the strict inequality
between successive coordinate sizes is used; the remaining reduction is
elementary linear algebra.  \leanmark{thinning}

\subsection{Removing one coordinate by a finite difference}

\begin{proposition}\label{prop:nondecomposition}
For every \(r\geq1\), there is no decomposition

\[
F_r=\sum_{i=0}^{r}f_i
\]

in which \(f_i\) is invariant under \(\sigma_r\) and almost constant in
coordinate \(i\).
\end{proposition}

\begin{proof}
We establish a reduction in degree and then prove the two-coordinate
base case separately.  Suppose that such a decomposition has been given.
For each \(y\in C_{r-1}\), choose an eventual value \(c_y\) and a finite
exceptional set \(N_y\subseteq D_r\) for the last-coordinate function
\(z\mapsto f_r(y,z)\).  Lemma~\ref{lem:uniform-thinning} supplies
\(a\in A_r\) avoiding all these exceptions.

For any function \(f:C_r\to\F\), define its two-bit difference by

\[
(\Delta_a f)(y)=f(y,(a,0))+f(y,(a,1)),
\qquad y\in C_{r-1}.
\]

For our chosen \(a\), both values of the last summand equal \(c_y\).
Consequently, \(\Delta_a f_r=0\).  In contrast, the product function
loses exactly its last factor:

\[
\begin{aligned}
(\Delta_a F_r)(y)
&=F_{r-1}(y)\bigl(s_{r-1}+(s_{r-1}+1)\bigr)\\
&=F_{r-1}(y).
\end{aligned}
\]

For \(i<r\), the function \(\Delta_a f_i\) remains almost constant in
coordinate \(i\).  Fix the other coordinates and take the union of the
finite exceptional sets for its two slices.  Outside this union, the sum
of their two eventual values is constant.

The difference also preserves invariance.  If \(f\) is invariant under
the simultaneous flip, then

\[
\begin{aligned}
(\Delta_a f)(\sigma_{r-1}y)
&=f(\sigma_{r-1}y,(a,0))+f(\sigma_{r-1}y,(a,1))\\
&=f(y,(a,1))+f(y,(a,0))\\
&=(\Delta_a f)(y).
\end{aligned}
\]

Here it matters that the two points have the same base coordinate:
the flip exchanges the two summands.  Thus, for \(r\geq2\), applying
\(\Delta_a\) gives

\[
F_{r-1}=\sum_{i=0}^{r-1}\Delta_a f_i,
\]

with all the required properties in one lower positive degree.  This
contradicts the induction hypothesis once the case \(r=1\) is known.

In this reduction, the largest coordinate and the summand controlled by
that coordinate disappear together.  Each surviving summand retains its
own almost-constancy condition, while the target function becomes the
same adjacent-bit product on the shorter configuration space.  The
operation sums only two values, so no infinite sums or convergence
questions enter the reduction.

For that base case, the reduction gives \(\Delta_a f_0=1\).  Retain the
single slice

\[
g(x)=f_0(x,(a,0)),
\qquad x\in D_0.
\]

It is almost constant because \(f_0\) has that property in coordinate
zero.  Invariance of \(f_0\), applied to \((x,(a,1))\), gives
\(g(\sigma_0x)=f_0(x,(a,1))\).  Hence

\[
g(x)+g(\sigma_0x)=1
\qquad\text{for every }x\in D_0.
\]

Choose a finite exceptional set \(N\subseteq D_0\) for \(g\), with
eventual value \(c\).  The set \(A_0\) is infinite, so it contains
\(a_0\) outside the first-coordinate projection of \(N\).  Both
\((a_0,0)\) and \((a_0,1)\) then lie outside \(N\).  Substituting the
first of these points into the displayed identity gives \(c+c=1\),
contradicting \(c+c=0\).  This proves the base case and the induction.\leanmark{difference}
\end{proof}

\begin{remark}
Two tempting shortcuts would lose the argument.  First, the single slice
\(g\) need not be invariant; the proof uses the relation between \(g\)
and its flipped value instead.  Second, one cannot continue a bare
nondecomposition induction down to degree zero.  The function \(F_0=1\)
is itself invariant and almost constant, so the corresponding
nondecomposition assertion is false.  The base case succeeds because it
retains the extra information that this constant function is the flip
difference of an almost-constant slice.
\end{remark}

\subsection{Returning to cohomology}

The top cochain in \eqref{eq:top-cochain} represents \(\eta_r^r\), by the
chain-lift calculation.  If this class were zero, the projective
resolution would express that cochain as a coboundary.
Lemma~\ref{lem:boundary-necessary} would then produce the decomposition
forbidden by Proposition~\ref{prop:nondecomposition}.  We conclude that

\[
\eta_r^r\neq0.
\]

The same resolution has length \(r\), so it also gives
\(H^{r+1}(U_r;\F)=0\), and therefore \(\eta_r^{r+1}=0\).  Combined with
the cardinality and weight of \(U_r\), these statements prove the
existence assertions of Theorem~\ref{thm:main}.  \leanmark{nonzero}
